\section{From accurate tracking to useful progress}\label{taskcontract}
An authoritative request says which goal is valid. An advisory estimate need not
change that goal. The executor assumes this distinction is correct; it does not
infer intent from noisy position packets.

Consider one new goal $g$, issued at $s$ and executed at $\tau\geq s\geq0$.
Every old reference event must finish by $s$, and its final target must be $a$.
Retain the entire old history and append $(\tau,g-a)$. The condition concerns
all old events, including future commitments. Zero velocity now does not imply
that no motion is queued.

\begin{theorem}[Isolated position task]
Let $h,J$ and the new event $(\tau,g-a)$ have exact rational encodings for the
checker. Assume the old rational history passes the exact jerk check, the stated
completed-tail condition holds, and
\[
\frac{2|g-a|}{h^3}\leq J,\qquad \tau+4h\leq D\leq H.
\]
Retain the plant, initial-error and budget assumptions of Theorem~\ref{tracking}
for the full appended history on $[0,H]$. Then the appended exact check passes.
For $s\leq t\leq H$, let $f=F((t-\tau)/h)$. The actual position satisfies
\begin{align*}
\min(a,g)-R&\leq x(t)\leq\max(a,g)+R,\\
|x(t)-g|&\leq |g-a|(1-f)+R.
\end{align*}
Between issue and execution, $|x(t)-a|\leq R$. At $\tau+2h$, if that time is in
the horizon, directed progress from $a$ toward $g$ is at least $|g-a|/2-R$.
For every $t\in[D,H]$, $|x(t)-g|\leq R$.
\end{theorem}

\textit{Proof idea.} Lean proves $0\leq F\leq1$, monotonicity, $F(2)=1/2$ and
the isolated jerk bound from the existing truncated-power primitive. Completed
old events reduce to $a$ after issue. Their jerk supports do not overlap the new
event's support, so the old guard and isolated bound imply the new exact guard.
The physical tracking theorem then supplies $|x-q|\leq R$. The displayed task
bounds follow from the range and midpoint properties. The physical trajectory
and its initial cube stay anchored at time zero; no state is reset at admission.

\textbf{Checked example.} For $a=0$, $g=20$ mm, $h=.2$ s, $s=.20$ s,
$\tau=.22$ s and $R=.2$ mm, the theorem gives at least $9.8$ mm of directed
progress at $.62$ s. Position is within $.2$ mm of the goal from $1.02$ s through
the chosen $2$ s horizon, including the $1.5$ s deadline. A separate checked
lower bound gives at least $9.8$ mm of goal error if the deadline is instead
$.62$ s. These are conditional analytic bounds, not measured robot performance.

\textbf{Scope.} Arrival is a position statement. It does not prove physical rest
or the benchmark's tight speed and acceleration conditions. The theorem assumes
the goal remains valid, bounds the declared trajectory, and requires a deadline
inside the covered horizon. It does not prove real communication timing,
arbitrary mid-motion replacement, or recognition of a small intentional move.

The source entry point is \texttt{checked\_isolated\_task}.
\par\path{task-contract-08/builder/TaskContract.lean}
The independent review and
negative controls are preserved beside it. This is a small composition result
that closes a task-level gap; it is not a new theory of intent or smoothing.
